% Source: https://artofproblemsolving.com/community/c3h1317780_counting_triangles
% Stable identity: aops-topic-1317780-q01
Write the subset as $a<b<c$; it is a triangle exactly when $a+b>c$. Fix the largest element $c$. We count pairs $a<b\le c-1$ with $a+b\ge c+1$. For each $b$, the condition is $c+1-b\le a\le b-1$, giving $\max(0,\,2b-c-1)$ choices. Summing over $b=1,\dots,c-1$ gives $1+3+5+\dots$ (for $c$ odd) or $2+4+6+\dots$ (for $c$ even), which in both cases equals $\left\lfloor\dfrac{(c-2)^2}{4}\right\rfloor$. Summing over $c=3,\dots,n$ and putting $k=c-2$ gives the formula.

For $n=10$:
\[
T(10)=\sum_{k=1}^{8}\left\lfloor\dfrac{k^2}{4}\right\rfloor=0+1+2+4+6+9+12+16=\boxed{50}.
\]
(The sum also has the closed form $\left\lfloor\dfrac{n(n-2)(2n-5)}{24}\right\rfloor$, which gives $\left\lfloor\dfrac{10\cdot8\cdot15}{24}\right\rfloor=50$.)

\begin{asy}
usepackage("amsmath");
size(5cm);
defaultpen(fontsize(10pt));
draw((0,0)--(8,0)); for(int i=0;i<=8;++i) draw((i,-0.1)--(i,0.1));
draw((0,0.4)--(3,0.4),blue+linewidth(1.5)); draw((0,0.8)--(4,0.8),blue+linewidth(1.5)); draw((0,1.2)--(6,1.2),red+linewidth(1.5));
label("$a$",(0,0.4),W,blue); label("$b$",(0,0.8),W,blue); label("$c$",(0,1.2),W,red); label("$a+b>c$",(7,1),E);
label("\mbox{}",(min(currentpicture,user=true)+max(currentpicture,user=true))/2,black); // keeps the text after the figure black when LaTeX typesets the labels inline
\end{asy}
